\begin{frame}{Методическая библиотека: \ProgLibraryBooks\ книг о подготовке и защите диссертации}
\centering
\begin{tikzpicture}[node distance=8mm,
  stage/.style={process,text width=40mm,minimum height=22mm}]
\node[stage] (ask) {Вопрос о тексте или докладе};
\node[stage,right=of ask] (search) {Поиск по книгам};
\node[stage,right=of search] (passage) {Отрывок с книгой и страницей};
\node[stage,right=of passage] (edit) {Правка слайда или текста};
\draw[flow] (ask)--(search);\draw[flow] (search)--(passage);\draw[flow] (passage)--(edit);
\end{tikzpicture}
\SlideFigureCaption{Использование методической библиотеки}
\SlideGap
\raggedright
\begin{itemize}
\item Авторы: \ProgLibraryAuthors.
\item Полных текстов --- \ProgLibraryFullText, фрагментов --- \ProgLibraryChunks; контрольные запросы найдены: \ProgLibraryHits\ из \ProgLibraryCases.
\item Советы книг не используются как научные доказательства. Первое применение --- презентация к защите; далее --- введение, автореферат и доклад.
\end{itemize}
\end{frame}
